\documentclass[11pt]{article}
\usepackage[margin=1in]{geometry}
\usepackage{amsmath,amssymb,amsthm,mathtools}
\usepackage{graphicx}
\usepackage[colorlinks=true,linkcolor=blue,citecolor=blue,urlcolor=blue]{hyperref}

\newtheorem{theorem}{Theorem}
\newtheorem{corollary}{Corollary}
\newtheorem{lemma}{Lemma}
\newtheorem{remark}{Remark}

\newcommand{\R}{\mathbb R}
\newcommand{\Z}{\mathbb Z}
\newcommand{\W}{\mathcal W}
\newcommand{\supp}{\operatorname{supp}}

\title{A No-Small-Alias Subcase of the Calderon Sum Formula for Wavelet Systems}
\author{agent\_lane\_02}
\date{2026-07-03}

\begin{document}
\maketitle

\section*{Source Problem}

The source paper is Ulrik Enstad and Jordy Timo van Velthoven,
\emph{On the Calderon sum formula for wavelet systems}, arXiv:2602.10766.
Its introduction records the following conjecture of Bownik--Lemvig and
Speegle: if an affine wavelet system associated with matrices
$A,P\in \operatorname{GL}(d,\R)$ and $\psi\in L^2(\R^d)$ is an orthonormal
basis, or more generally a Parseval frame, then
\[
  \sum_{j\in\Z} |\widehat\psi((A^t)^j\xi)|^2=|\det P|
  \quad\text{for a.e. }\xi\in\R^d.
\]
The source proves this under the additional regularity hypothesis
$\psi\in\mathcal B_\pi$.  The result below gives a different subcase which
does not assume $\psi\in\mathcal B_\pi$.

\begin{center}
\includegraphics[width=0.95\linewidth]{figures/open_problem_crop.png}
\end{center}

\section*{Statement}

We use the quasi-regular representation notation of Section 4 of the source:
\[
  \pi(x,A^j)f(t)=|\det A|^{-j/2} f(A^{-j}(t-x)).
\]
Thus
\[
  \widehat{\pi(A^jPk,A^j)\psi}(\xi)
  =|\det A|^{j/2} e^{-2\pi i\,\xi\cdot A^jPk}
    \widehat\psi((A^t)^j\xi).
\]

\begin{theorem}[No-small-alias subcase]\label{thm:main}
Let $A,P\in\operatorname{GL}(d,\R)$ and $\psi\in L^2(\R^d)$.  Assume
\[
  \delta(A,P):=
  \inf_{\substack{j\in\Z\\ m\in\Z^d\setminus\{0\}}}
  \|(A^jP)^{-t}m\|>0.
\]
If
\[
  \W(\psi,A,P):=
  \{\pi(A^jPk,A^j)\psi: j\in\Z,\ k\in\Z^d\}
\]
is a Parseval frame for $L^2(\R^d)$, then
\[
  \sum_{j\in\Z} |\widehat\psi((A^t)^j\xi)|^2=|\det P|
  \quad\text{for a.e. }\xi\in\R^d.
\]
More generally, if $\W(\psi,A,P)$ has frame bounds $0<C_1\le C_2<\infty$,
then
\[
  C_1|\det P|
  \le
  \sum_{j\in\Z} |\widehat\psi((A^t)^j\xi)|^2
  \le
  C_2|\det P|
  \quad\text{for a.e. }\xi\in\R^d.
\]
\end{theorem}

\section*{Proof Intuition}

The obstruction in the full conjecture is aliasing.  At scale $j$, summing the
translation coefficients over $k\in\Z^d$ periodizes in frequency along the dual
lattice $(A^jP)^{-t}\Z^d$.  If dual lattice vectors from different scales can
be arbitrarily small, then no tiny Fourier neighborhood is alias-free at every
scale; this is exactly where the source paper needs the coefficient-tail
condition $\psi\in\mathcal B_\pi$.

Under the no-small-alias hypothesis, we can choose Fourier supports with
diameter below $\delta(A,P)$.  On such supports every scale has only its
diagonal alias.  The Parseval frame identity then reads locally as integration
against the Calderon sum.  Since small cubes cover frequency space, the local
identity forces the Calderon sum formula almost everywhere.

\section*{Proof}

We first record the elementary lattice calculation used at each scale.

\begin{lemma}[Alias-free lattice summation]\label{lem:lattice}
Let $B\in\operatorname{GL}(d,\R)$, let $E\subset\R^d$ be measurable, and assume
\[
  (E-E)\cap \bigl(B^{-t}\Z^d\setminus\{0\}\bigr)=\emptyset.
\]
If $G\in L^2(\R^d)$ is supported in $E$, then
\[
  \sum_{k\in\Z^d}
  \left|\int_{\R^d} G(\xi)e^{2\pi i\,\xi\cdot Bk}\,d\xi\right|^2
  =
  |\det B|^{-1}\int_E |G(\xi)|^2\,d\xi .
\]
\end{lemma}

\begin{proof}
Let $Q$ be a measurable fundamental domain for the dual lattice
$B^{-t}\Z^d$.  Periodize
\[
  G_B(\eta)=\sum_{\ell\in B^{-t}\Z^d}G(\eta+\ell),
  \qquad \eta\in Q.
\]
The assumption on $E-E$ implies that at most one summand is nonzero for a.e.
$\eta\in Q$, hence
\[
  \int_Q |G_B(\eta)|^2\,d\eta=\int_E |G(\xi)|^2\,d\xi.
\]
Moreover
\[
  \int_{\R^d}G(\xi)e^{2\pi i\,\xi\cdot Bk}\,d\xi
  =
  \int_Q G_B(\eta)e^{2\pi i\,\eta\cdot Bk}\,d\eta .
\]
The functions $|\det B|^{1/2}e^{2\pi i\,\eta\cdot Bk}$, $k\in\Z^d$, form an
orthonormal basis of $L^2(Q)$.  Parseval's identity on $Q$ gives the claim.
\end{proof}

We now prove the theorem.  Put
\[
  h_j(\xi):=|\det A|^{j/2}\widehat\psi((A^t)^j\xi),
  \qquad B_j:=A^jP .
\]
Then
\[
  \widehat{\pi(A^jPk,A^j)\psi}(\xi)
  =e^{-2\pi i\,\xi\cdot B_jk}h_j(\xi).
\]

Let $E\subset\R^d$ have diameter strictly smaller than $\delta(A,P)$, and let
$f\in L^2(\R^d)$ be such that $\widehat f$ is bounded and supported in $E$.
For every $j$, the no-small-alias condition gives
\[
  (E-E)\cap \bigl(B_j^{-t}\Z^d\setminus\{0\}\bigr)=\emptyset.
\]
Applying Lemma~\ref{lem:lattice} with
$G(\xi)=\widehat f(\xi)\overline{h_j(\xi)}$ yields
\[
\begin{aligned}
  \sum_{k\in\Z^d}
  |\langle f,\pi(A^jPk,A^j)\psi\rangle|^2
  &=
  |\det B_j|^{-1}\int_E |\widehat f(\xi)|^2 |h_j(\xi)|^2\,d\xi  \\
  &=
  |\det P|^{-1}\int_E |\widehat f(\xi)|^2
      |\widehat\psi((A^t)^j\xi)|^2\,d\xi .
\end{aligned}
\]
Summing over $j$ and using Tonelli's theorem, the frame inequalities for
$\W(\psi,A,P)$ give
\[
  C_1\int_E|\widehat f(\xi)|^2\,d\xi
  \le
  |\det P|^{-1}\int_E|\widehat f(\xi)|^2
       \sum_{j\in\Z}|\widehat\psi((A^t)^j\xi)|^2\,d\xi
  \le
  C_2\int_E|\widehat f(\xi)|^2\,d\xi .
\]
Since this holds for all bounded $\widehat f$ supported in $E$, it follows
that
\[
  C_1|\det P|
  \le
  \sum_{j\in\Z}|\widehat\psi((A^t)^j\xi)|^2
  \le
  C_2|\det P|
  \quad\text{for a.e. }\xi\in E.
\]
Finally, $\R^d$ is a countable union of cubes of diameter less than
$\delta(A,P)$.  Applying the preceding conclusion to each cube proves the
global a.e. statement.  In the Parseval case $C_1=C_2=1$, this is precisely
the Calderon sum formula.
\qed

\section*{Examples And Scope}

The condition $\delta(A,P)>0$ is an arithmetic separation condition on the
reciprocal translation lattices.  It holds, for example, whenever the lattices
$(A^jP)^{-t}\Z^d$ range over a uniformly discrete family.  A concrete class is
given by lattice automorphism regimes: after conjugating by $P$, if $A^{-t}$
acts as an automorphism of a fixed lattice, then all reciprocal lattices are
the same lattice and $\delta(A,P)>0$.

The result does not cover the usual expansive case.  If $A$ is expansive, then
nonzero vectors of $(A^jP)^{-t}\Z^d$ typically approach zero as $j\to+\infty$,
so the alias-free localization argument cannot start.  That is consistent with
the source paper's use of a different mechanism, namely the coefficient-space
condition $\psi\in\mathcal B_\pi$.

\section*{Novelty And Verification Notes}

Local run indexes were searched for arXiv:2602.10766 and for Calderon/wavelet
keywords; no existing exact packet was found.  Web searches for the exact
Calderon-sum conjecture and Bownik--Lemvig problem located the source arXiv
paper but no later resolution.  This packet claims only the no-small-alias
subcase, not the full conjecture.

The proof uses only Plancherel/Parseval on the dual lattice and the frame
identity.  No computational verification is involved.

\section*{References}

\begin{enumerate}
\item U. Enstad and J. T. van Velthoven,
\emph{On the Calderon sum formula for wavelet systems}, arXiv:2602.10766.
\item M. Bownik and J. Lemvig, references cited in the source as the origin of
the Calderon-sum conjecture/open problem.
\item D. Speegle, reference cited in the source as an earlier implicit
formulation of the same problem.
\end{enumerate}

\section*{Human Review Recommendation}

Recommended review focus: check the normalization convention between the
source's displayed affine system and its Section 4 representation notation,
and verify the determinant factor in Lemma~\ref{lem:lattice}.  The mathematical
mechanism is otherwise compact: the no-small-alias condition removes every
off-diagonal alias term.

\end{document}
